\documentclass{article}
\usepackage{amsmath,amssymb,amsthm}
\usepackage{algorithm}
\usepackage[noend]{algpseudocode}
\usepackage{enumitem}
\usepackage{hyperref}
\usepackage{geometry}
\geometry{margin=1in}
\newtheorem{theorem}{Theorem}
\newtheorem{lemma}{Lemma}
\newtheorem{assumption}{Assumption}
\newcommand{\proj}[2]{\Pi_{#1}\!\left(#2\right)}
\newcommand{\gap}{\mathcal{G}}
\newcommand{\R}{\mathbb{R}}

\begin{document}

\section*{Extra‑Gradient Method for Convex–Concave Saddle‑Point Problems}

\section*{Mathematical Formulation}
Consider the saddle‑point problem  

\[
\min_{x\in X}\;\max_{y\in Y}\;L(x,y),
\]

where  

\begin{itemize}
  \item $X\subseteq\R^{n}$ and $Y\subseteq\R^{m}$ are closed, convex sets,
  \item $L:\R^{n}\times\R^{m}\rightarrow\R$ is convex in $x$ and concave in $y$,
  \item $L$ is continuously differentiable and its gradient is $L$‑Lipschitz:
  \[
  \|\nabla L(u)-\nabla L(v)\| \le L\|u-v\|\qquad\forall u,v\in\R^{n+m}.
  \]
\end{itemize}

Define the monotone operator  

\[
F(z)\;:=\;
\begin{pmatrix}
\nabla_{x}L(x,y)\\[2mm]
-\nabla_{y}L(x,y)
\end{pmatrix},
\qquad z:=(x,y).
\]

The extra‑gradient (Korpelevich) scheme with stepsize $\eta>0$ reads  

\[
\begin{aligned}
z_{k+1/2}&:=\proj{X\times Y}{\,z_{k}-\eta F(z_{k})\,},\\[1mm]
z_{k+1}&:=\proj{X\times Y}{\,z_{k}-\eta F(z_{k+1/2})\,}.
\end{aligned}
\tag{EG}
\]

When $X=\R^{n},Y=\R^{m}$ the projections are identities and the updates simplify to  

\[
\begin{aligned}
x_{k+1/2}&=x_{k}-\eta \nabla_{x}L(x_{k},y_{k}), &
y_{k+1/2}&=y_{k}+ \eta \nabla_{y}L(x_{k},y_{k}),\\[1mm]
x_{k+1}&=x_{k}-\eta \nabla_{x}L(x_{k+1/2},y_{k+1/2}), &
y_{k+1}&=y_{k}+ \eta \nabla_{y}L(x_{k+1/2},y_{k+1/2}).
\end{aligned}
\tag{1}
\]

\section*{Pseudocode}
\begin{algorithm}[H]
\caption{Extra‑Gradient for $\min_{x}\max_{y}L(x,y)$}
\begin{algorithmic}[1]
\State \textbf{Input:} stepsize $\eta>0$, initial $(x_{0},y_{0})\in X\times Y$, max iter $T$.
\For{$k=0,1,\dots,T-1$}
    \State Compute gradients at current point:
    \[
    g^{x}_{k}\gets \nabla_{x}L(x_{k},y_{k}),\qquad 
    g^{y}_{k}\gets \nabla_{y}L(x_{k},y_{k})
    \]
    \State Extrapolation:
    \[
    \tilde{x}\gets \proj{X}{\,x_{k}-\eta g^{x}_{k}\,},\qquad
    \tilde{y}\gets \proj{Y}{\,y_{k}+\eta g^{y}_{k}\,}
    \]
    \State Gradients at extrapolated point:
    \[
    \tilde{g}^{x}_{k}\gets \nabla_{x}L(\tilde{x},\tilde{y}),\qquad
    \tilde{g}^{y}_{k}\gets \nabla_{y}L(\tilde{x},\tilde{y})
    \]
    \State Update:
    \[
    x_{k+1}\gets \proj{X}{\,x_{k}-\eta \tilde{g}^{x}_{k}\,},\qquad
    y_{k+1}\gets \proj{Y}{\,y_{k}+\eta \tilde{g}^{y}_{k}\,}
    \]
\EndFor
\State \textbf{Output:} $(x_{T},y_{T})$ or the ergodic average
$\displaystyle \bar z_{T}=\frac1T\sum_{k=0}^{T-1}(x_{k},y_{k})$.
\end{algorithmic}
\end{algorithm}

\section*{Dry Run Example}
We illustrate (EG) on the simple convex minimisation problem  

\[
\min_{w\in\R} \; f(w):=(w-3)^{2},
\]

which can be written as a saddle‑point with $y$ absent.  
Here $L(w)=f(w)$, $\nabla f(w)=2(w-3)$ and the projection is the identity.

\begin{center}
\begin{tabular}{c|c|c|c}
iteration $k$ & $w_{k}$ & $\nabla f(w_{k})$ & $w_{k+1}$ \\\hline
0 & $0$ & $2(0-3)=-6$ & $w_{1}=0-0.1\bigl[\,2\bigl(0-0.1\cdot(-6)-3\bigr)\bigr]=0.12$\\[2mm]
1 & $0.12$ & $2(0.12-3)=-5.76$ &
$w_{2}=0.12-0.1\bigl[\,2\bigl(0.12-0.1\cdot(-5.76)-3\bigr)\bigr]=0.2304$\\[2mm]
2 & $0.2304$ & $2(0.2304-3)=-5.5392$ &
$w_{3}=0.2304-0.1\bigl[\,2\bigl(0.2304-0.1\cdot(-5.5392)-3\bigr)\bigr]=0.3305\;(\text{rounded})$
\end{tabular}
\end{center}

The objective values decrease: $f(0)=9$, $f(0.12)=8.2944$, $f(0.2304)=7.6176$, $f(0.3305)\approx6.979$,
illustrating the benefit of the extrapolation step.

\section*{Assumptions}
\begin{assumption}[Convex–Concave Structure]\label{as:cc}
$L(\cdot ,y)$ is convex for every $y\in Y$ and $L(x,\cdot )$ is concave for every $x\in X$.
\end{assumption}
\begin{assumption}[Lipschitz Gradient]\label{as:lips}
There exists $L>0$ such that for all $z_{1},z_{2}\in X\times Y$,
\[
\|F(z_{1})-F(z_{2})\|\le L\|z_{1}-z_{2}\|.
\]
\end{assumption}
\begin{assumption}[Monotonicity]\label{as:mono}
$F$ is monotone, i.e.
\[
\langle F(z_{1})-F(z_{2}),\,z_{1}-z_{2}\rangle\ge0,\qquad\forall z_{1},z_{2}\in X\times Y.
\]
\end{assumption}
\begin{assumption}[Stepsize]\label{as:step}
The stepsize satisfies $0<\eta\le \frac{1}{2L}$.
\end{assumption}

\section*{Main Convergence Theorem}
\begin{theorem}[Ergodic $O(1/T)$ Rate]\label{thm:rate}
Let $\{z_{k}\}_{k\ge0}$ be generated by \eqref{EG} under Assumptions \ref{as:cc}–\ref{as:step},
and let $z^{*}=(x^{*},y^{*})$ be any saddle point.  
Define the averaged iterate  

\[
\bar z_{T}:=\frac1T\sum_{k=0}^{T-1}z_{k}.
\]

Then the primal–dual gap satisfies  

\[
\gap(\bar z_{T})\;:=\;\max_{y\in Y}L(\bar x_{T},y)-\min_{x\in X}L(x,\bar y_{T})
\;\le\;\frac{\|z_{0}-z^{*}\|^{2}}{2\eta T}.
\tag{2}
\]

Consequently $\gap(\bar z_{T})=O(1/T)$.
\end{theorem}

\section*{Theoretical Properties}
\begin{itemize}
\item \textbf{Stability.} The condition $\eta\le 1/(2L)$ guarantees that the operator
$T_{\eta}(z):=\proj{X\times Y}{z-\eta F(z)}$ is non‑expansive, which is crucial for the
contraction argument in the proof.
\item \textbf{Hyperparameter Sensitivity.} If $\eta$ is chosen larger than $1/(2L)$,
the monotonicity inequality may be violated and the bound in \eqref{2} can become
negative, i.e. the method may diverge. Hence the bound is tight in the sense
that $1/(2L)$ is the largest admissible constant for the standard analysis.
\item \textbf{Comparison.} Classical gradient descent on the primal problem alone
yields an $O(1/T)$ rate for convex smooth functions but requires solving the
maximisation in $y$ exactly at each iteration. The extra‑gradient method avoids
that inner optimisation while preserving the same ergodic rate.
\end{itemize}

\section*{Preliminaries (Definitions and Lemmas)}
\begin{lemma}[Non‑expansiveness of Projection]\label{lem:proj}
For any closed convex set $C\subseteq\R^{d}$ and any $u,v\in\R^{d}$,
\[
\|\proj{C}{u}-\proj{C}{v}\|\le\|u-v\|.
\]
\end{lemma}
\begin{proof}
Standard property of Euclidean projections (see e.g., \cite{Bauschke2011}).
\end{proof}

\begin{lemma}[Descent Lemma for Monotone Lipschitz Operators]\label{lem:descent}
Let $F$ satisfy Assumptions \ref{as:lips} and \ref{as:mono}. For any $z\in X\times Y$ and stepsize
$\eta\le 1/(2L)$,
\[
\| \proj{X\times Y}{z-\eta F(z)}-z^{*}\|^{2}
\le \|z-z^{*}\|^{2}-\eta\langle F(z),z-z^{*}\rangle .
\tag{3}
\]
\end{lemma}
\begin{proof}
Let $p:=\proj{X\times Y}{z-\eta F(z)}$. By Lemma~\ref{lem:proj} and the definition of $p$,
\[
\|p-z^{*}\|^{2}
= \|z-\eta F(z)-z^{*}\|^{2}
- \|z-\eta F(z)-p\|^{2}
\le \|z-\eta F(z)-z^{*}\|^{2}.
\]
Expanding the RHS,
\[
\|z-z^{*}\|^{2}-2\eta\langle F(z),z-z^{*}\rangle+\eta^{2}\|F(z)\|^{2}.
\]
Using $\|F(z)\|\le L\|z-z^{*}\|$ (Lipschitzness with $F(z^{*})=0$) and $\eta\le 1/(2L)$,
\[
\eta^{2}\|F(z)\|^{2}\le \eta^{2}L^{2}\|z-z^{*}\|^{2}\le \frac{\eta}{2}\|F(z)\|^{2}
\le \eta\langle F(z),z-z^{*}\rangle,
\]
where the last inequality follows from monotonicity $ \langle F(z)-F(z^{*}),z-z^{*}\rangle\ge0$.
Subtracting the term $\eta\langle F(z),z-z^{*}\rangle$ yields \eqref{3}.
\end{proof}

\section*{Main Proof (Step‑by‑Step Derivation)}
We prove Theorem~\ref{thm:rate}.  Throughout we denote $z_{k}=(x_{k},y_{k})$ and
$z^{*}=(x^{*},y^{*})$ a saddle point (hence $F(z^{*})=0$).

\begin{enumerate}[label=[Step \arabic*], leftmargin=*]
\item \textbf{Apply Lemma~\ref{lem:descent} to the extrapolation step.}  
For $z_{k}$ we have
\[
\|z_{k+1/2}-z^{*}\|^{2}\le\|z_{k}-z^{*}\|^{2}-\eta\langle F(z_{k}),z_{k}-z^{*}\rangle .
\tag{4}
\]

\item \textbf{Apply Lemma~\ref{lem:descent} to the correction step.}  
Using $z_{k+1/2}$ as the base point,
\[
\|z_{k+1}-z^{*}\|^{2}\le\|z_{k}-z^{*}\|^{2}
-\eta\langle F(z_{k+1/2}),z_{k}-z^{*}\rangle .
\tag{5}
\]

\item \textbf{Subtract \eqref{4} from \eqref{5}.}  
\[
\|z_{k+1}-z^{*}\|^{2}-\|z_{k+1/2}-z^{*}\|^{2}
\le -\eta\bigl[\langle F(z_{k+1/2})-F(z_{k}),\,z_{k}-z^{*}\rangle\bigr].
\tag{6}
\]

\item \textbf{Bound the inner product using Lipschitzness.}  
By Cauchy–Schwarz and Assumption~\ref{as:lips},
\[
\bigl|\langle F(z_{k+1/2})-F(z_{k}),\,z_{k}-z^{*}\rangle\bigr|
\le \|F(z_{k+1/2})-F(z_{k})\|\,\|z_{k}-z^{*}\|
\le L\|z_{k+1/2}-z_{k}\|\,\|z_{k}-z^{*}\|.
\tag{7}
\]

\item \textbf{Relate $\|z_{k+1/2}-z_{k}\|$ to the gradient at $z_{k}$.}  
From the definition of $z_{k+1/2}$,
\[
z_{k+1/2}= \proj{X\times Y}{z_{k}-\eta F(z_{k})},
\]
hence by Lemma~\ref{lem:proj},
\[
\|z_{k+1/2}-z_{k}\|\le \eta\|F(z_{k})\|.
\tag{8}
\]

\item \textbf{Insert (8) into (7) and then into (6).}  
\[
\|z_{k+1}-z^{*}\|^{2}-\|z_{k+1/2}-z^{*}\|^{2}
\le \eta L\,\eta\|F(z_{k})\|\|z_{k}-z^{*}\|
= \eta^{2}L\|F(z_{k})\|\|z_{k}-z^{*}\|.
\tag{9}
\]

\item \textbf{Upper‑bound $\|F(z_{k})\|$ by monotonicity.}  
Monotonicity with $z^{*}$ gives
\[
\langle F(z_{k}),z_{k}-z^{*}\rangle\ge0
\quad\Longrightarrow\quad
\|F(z_{k})\|\,\|z_{k}-z^{*}\|\ge\langle F(z_{k}),z_{k}-z^{*}\rangle .
\tag{10}
\]

\item \textbf{Combine (4) and (9).}  
From (4),
\[
\|z_{k+1/2}-z^{*}\|^{2}\le\|z_{k}-z^{*}\|^{2}-\eta\langle F(z_{k}),z_{k}-z^{*}\rangle .
\tag{11}
\]
Adding (9) to (11) yields
\[
\|z_{k+1}-z^{*}\|^{2}
\le\|z_{k}-z^{*}\|^{2}
-\eta\langle F(z_{k}),z_{k}-z^{*}\rangle
+\eta^{2}L\|F(z_{k})\|\|z_{k}-z^{*}\|.
\tag{12}
\]

\item \textbf{Use $\eta\le 1/(2L)$ to absorb the last term.}  
From (10) and $\eta L\le \tfrac12$,
\[
\eta^{2}L\|F(z_{k})\|\|z_{k}-z^{*}\|
\le \frac{\eta}{2}\langle F(z_{k}),z_{k}-z^{*}\rangle .
\tag{13}
\]
Plugging (13) into (12) gives a clean descent inequality
\[
\|z_{k+1}-z^{*}\|^{2}
\le\|z_{k}-z^{*}\|^{2}
-\frac{\eta}{2}\langle F(z_{k}),z_{k}-z^{*}\rangle .
\tag{14}
\]

\item \textbf{Relate the inner product to the primal–dual gap.}  
For monotone $F$, the following variational inequality holds for any $z$:
\[
\langle F(z),z-z^{*}\rangle \ge \gap(z).
\tag{15}
\]
(Proof: by definition of saddle point and convex–concave structure.)

\item \textbf{Insert (15) into (14).}  
\[
\|z_{k+1}-z^{*}\|^{2}
\le\|z_{k}-z^{*}\|^{2}
-\frac{\eta}{2}\,\gap(z_{k}).
\tag{16}
\]

\item \textbf{Telescoping sum.}  
Sum \eqref{16} for $k=0,\dots,T-1$:
\[
\|z_{T}-z^{*}\|^{2}
\le\|z_{0}-z^{*}\|^{2}
-\frac{\eta}{2}\sum_{k=0}^{T-1}\gap(z_{k}).
\tag{17}
\]

\item \textbf{Discard the non‑negative term $\|z_{T}-z^{*}\|^{2}$ and divide.}
\[
\frac{1}{T}\sum_{k=0}^{T-1}\gap(z_{k})
\le \frac{2\|z_{0}-z^{*}\|^{2}}{\eta T}.
\tag{18}
\]

\item \textbf{Use convexity of the gap to pass to the average iterate.}  
Since $\gap(\cdot)$ is convex,
\[
\gap(\bar z_{T})\le\frac{1}{T}\sum_{k=0}^{T-1}\gap(z_{k}).
\tag{19}
\]

\item \textbf{Combine (18) and (19).}  
\[
\boxed{\;\gap(\bar z_{T})\le\frac{\|z_{0}-z^{*}\|^{2}}{2\eta T}\;}
\]
which is exactly the bound \eqref{2}.  ∎
\end{enumerate}

\section*{Rate Derivation (With Constants)}
The constant in \eqref{2} is $C:=\|z_{0}-z^{*}\|^{2}/(2\eta)$.  
If $X$ and $Y$ are bounded with diameters $D_{x},D_{y}$, then $\|z_{0}-z^{*}\|\le\sqrt{D_{x}^{2}+D_{y}^{2}}$,
so a fully explicit rate reads  

\[
\gap(\bar z_{T})\le
\frac{D_{x}^{2}+D_{y}^{2}}{4\eta T},\qquad
\eta\le\frac{1}{2L}.
\]

Choosing the largest admissible stepsize $\eta=1/(2L)$ yields  

\[
\gap(\bar z_{T})\le\frac{L\,(D_{x}^{2}+D_{y}^{2})}{2T}=O\!\left(\frac{1}{T}\right).
\]

\section*{Special Cases}
\begin{itemize}
\item \textbf{Unconstrained case $X=\R^{n},Y=\R^{m}$.}  
Projections disappear and the proof simplifies because the iterates remain in the whole space.
\item \textbf{Pure minimisation ($y$ absent).}  
The algorithm reduces to the classic \emph{extragradient descent} for smooth convex functions, and the bound \eqref{2} becomes the standard $O(1/T)$ rate for the function value gap.
\item \textbf{Strongly monotone $F$.}  
If $F$ is $\mu$‑strongly monotone, a similar derivation yields a linear rate
$\|z_{k}-z^{*}\|\le (1-\eta\mu)^{k}\|z_{0}-z^{*}\|$.
\end{itemize}

\begin{thebibliography}{9}
\bibitem{Bauschke2011}
H.~H. Bauschke and P.~L. Combettes,
\emph{Convex Analysis and Monotone Operator Theory in Hilbert Spaces},
Springer, 2011.
\end{thebibliography}

\end{document}